\documentclass[FileMain.tex]{subfiles}
\gdef\monhoc{Khoa học tự nhiên 6}
\gdef\ngaykt{23/05/2026}
\gdef\nh{2025 - 2026}
\gdef\thoigian{45}
\gdef\made{523}
\begin{document}
\section[Kiểm tra định kỳ KHTN 6 - Tách chất ra khỏi hỗn hợp]{Kiểm tra định kỳ KHTN 6 }

%%%==============Phần trắc nghiệm nhiều lựa chọn==============%%%
\subsection{Bài tập trắc nghiệm nhiều lựa chọn}\textit{\large Thí sinh trả lời từ câu 1 đến câu 4. Mỗi câu thí sinh chỉ chọn một phương án}
\Opensolutionfile{ansex}[Ans/LGEX-KTDK_KHTN_6_TachChat_MADE523]
\Opensolutionfile{ans}[Ans/Ans-KTDK_KHTN_6_TachChat_MADE523]

%%%%%============EX_1================%%%%%%
\begin{ex}%[6K3N1-1]
	Để tách muối ăn ra khỏi hỗn hợp muối ăn và cát, ta có thể sử dụng phương pháp nào sau đây?
	\choice
	{Dùng nam châm}
	{\True Hoà tan vào nước, lọc rồi cô cạn}
	{Dùng kính lúp để nhặt từng hạt muối}
	{Nung nóng hỗn hợp ở nhiệt độ cao}
	\loigiai{
		Muối ăn tan được trong nước còn cát không tan. Ta hoà tan hỗn hợp vào nước, lọc bỏ cát (phần không tan), rồi cô cạn nước lọc để thu được muối ăn tinh khiết.
	}
\end{ex}

%%%%%============EX_2================%%%%%%
\begin{ex}%[6K3N1-2]
	Phương pháp nào sau đây dùng để tách sắt ra khỏi hỗn hợp bột sắt và bột lưu huỳnh?
	\choice
	{Cô cạn}
	{Lọc}
	{Chiết}
	{\True Dùng nam châm}
	\loigiai{
		Sắt là kim loại có tính từ, bị nam châm hút. Lưu huỳnh không bị nam châm hút. Khi đưa nam châm vào hỗn hợp, sắt sẽ bị hút ra, còn lại lưu huỳnh. Đây là phương pháp tách dựa vào tính chất vật lý khác nhau của các chất.
	}
\end{ex}

%%%%%============EX_3================%%%%%%
\begin{ex}%[6K3H1-1]
	Để tách dầu ăn ra khỏi hỗn hợp dầu ăn và nước, ta sử dụng phương pháp nào?
	\choice
	{Lọc bằng giấy lọc}
	{Cô cạn}
	{\True Chiết}
	{Dùng nam châm}
	\loigiai{
		Dầu ăn không tan trong nước và nhẹ hơn nước nên nổi lên trên tạo thành hai lớp chất lỏng. Ta dùng phương pháp chiết: cho hỗn hợp vào phễu chiết, mở khoá phễu để nước (lớp dưới) chảy ra trước, giữ lại dầu ăn (lớp trên).
	}
\end{ex}

%%%%%============EX_4================%%%%%%
\begin{ex}%[6K3H1-2]
	Nước cất thu được bằng phương pháp chưng cất là
	\choice
	{hỗn hợp đồng nhất}
	{hỗn hợp không đồng nhất}
	{huyền phù}
	{\True chất tinh khiết}
	\loigiai{
		Chưng cất là phương pháp tách chất dựa vào sự khác nhau về nhiệt độ sôi. Khi đun sôi nước tự nhiên, hơi nước bốc lên được ngưng tụ thành nước cất. Nước cất chỉ chứa duy nhất nước ($H_2O$), không có tạp chất nên là chất tinh khiết.
	}
\end{ex}

\Closesolutionfile{ans}
\Closesolutionfile{ansex}
%\bangdapan{Ans-KTDK_KHTN_6_TachChat_MADE523}

%%%==============Phần trắc nghiệm đúng sai==============%%%
\subsection{Trắc nghiệm đúng sai}\textit{\large Thí sinh trả lời từ câu 1 đến câu 2. Trong mỗi ý a), b), c), d) ở mỗi câu thí sinh chọn đúng hoặc sai}
\Opensolutionfile{ansex}[Ans/LGTF-KTDK_KHTN_6_TachChat_MADE523]
\Opensolutionfile{ansbook}[Ansbook/AnsTF-KTDK_KHTN_6_TachChat_MADE523]
\setcounter{ex}{0}

%%%%%============TF_1================%%%%%%
\begin{ex}%[6K3H2-1]
	Trong phòng thí nghiệm, để thu được nước tinh khiết từ nước muối, bạn An đã tiến hành chưng cất. Đánh giá tính đúng/sai của các phát biểu sau:
	\choiceTF
	{\True Nước muối là hỗn hợp đồng nhất gồm nước và muối ăn}
	{\True Khi đun sôi nước muối, nước bay hơi trước vì có nhiệt độ sôi thấp hơn muối}
	{Sau khi chưng cất, trong bình hứng sẽ thu được muối ăn tinh khiết}
	{\True Phương pháp chưng cất dựa vào sự khác nhau về nhiệt độ sôi của các chất}
	\loigiai{
		\begin{itemchoice}[T1,T2,F3,T4]
			\itemch Nước muối là dung dịch gồm chất tan (muối ăn $NaCl$) tan trong dung môi (nước $H_2O$), là hỗn hợp đồng nhất. Đúng.
			\itemch Nhiệt độ sôi của nước là $100\,^{\circ}\text{C}$, thấp hơn rất nhiều so với nhiệt độ sôi của muối ăn ($1413\,^{\circ}\text{C}$). Khi đun sôi, nước sẽ bay hơi trước. Đúng.
			\itemch Sau khi chưng cất, nước bay hơi rồi ngưng tụ lại trong bình hứng. Vậy bình hứng thu được nước cất (chất tinh khiết), không phải muối. Muối ăn còn lại trong bình đun. Sai.
			\itemch Chưng cất tách các chất dựa vào sự khác nhau về nhiệt độ sôi: chất có nhiệt độ sôi thấp hơn sẽ bay hơi trước, được ngưng tụ và thu lại. Đúng.
		\end{itemchoice}
	}
\end{ex}

%%%%%============TF_2================%%%%%%
\begin{ex}%[6K3V2-1]
	Bạn Minh có một hỗn hợp gồm bột sắt, muối ăn và cát. Bạn ấy muốn tách riêng từng chất. Đánh giá tính đúng/sai của các phát biểu sau:
	\choiceTF
	{\True Bước đầu tiên nên dùng nam châm để tách bột sắt ra khỏi hỗn hợp}
	{Có thể dùng phương pháp lọc để tách muối ăn ra khỏi cát mà không cần hoà tan vào nước}
	{\True Sau khi tách sắt, hoà tan hỗn hợp còn lại vào nước thì muối tan, cát không tan}
	{Để thu muối ăn từ nước muối ta dùng phương pháp chiết}
	\loigiai{
		\begin{itemchoice}[T1,F2,T3,F4]
			\itemch Sắt có tính từ, bị nam châm hút. Dùng nam châm là cách đơn giản và hiệu quả nhất để tách sắt ra trước. Đúng.
			\itemch Muối ăn và cát đều là chất rắn dạng hạt. Nếu không hoà tan vào nước thì không thể dùng giấy lọc để tách riêng chúng, vì cả hai đều không lọt qua lỗ giấy lọc. Cần hoà tan vào nước trước rồi mới lọc. Sai.
			\itemch Muối ăn tan tốt trong nước tạo dung dịch, cát không tan trong nước. Khi lọc, cát nằm lại trên giấy lọc, nước muối chảy qua. Đúng.
			\itemch Để thu muối ăn từ nước muối, ta dùng phương pháp cô cạn (bay hơi nước), không dùng chiết. Chiết dùng để tách hai chất lỏng không hoà tan vào nhau. Sai.
		\end{itemchoice}
	}
\end{ex}

\Closesolutionfile{ansbook}
\Closesolutionfile{ansex}
%\bangdapanTF{AnsTF-KTDK_KHTN_6_TachChat_MADE523}

%%%==============Phần bài tập trả lời ngắn==============%%%
\subsection{Bài tập trả lời ngắn}\textit{\large Thí sinh trả lời từ câu 1 đến câu 2}
\Opensolutionfile{ansex}[Ans/LGSA-KTDK_KHTN_6_TachChat_MADE523]
\Opensolutionfile{ansexh}[Ans/AnsSA-KTDK_KHTN_6_TachChat_MADE523]
\setcounter{ex}{0}

%%%%%============SA_1================%%%%%%
\begin{ex}%[6K3H3-1]
	Người ta cô cạn $200$ gam dung dịch nước muối có nồng độ $5\%$ để thu được muối ăn tinh khiết. Hỏi khối lượng muối ăn thu được là bao nhiêu gam?
	\shortans{$10$}
	\loigiai{
		Khối lượng muối ăn thu được bằng khối lượng chất tan trong dung dịch.
		\\
		Nồng độ phần trăm của dung dịch là $5\%$, nghĩa là trong $100$ gam dung dịch có $5$ gam muối ăn.
		\\
		Khối lượng muối ăn thu được:
		\[ m_{\text{muối}} = \dfrac{5}{100} \times 200 = 10 \text{ (gam)} \]
		Vậy sau khi cô cạn, ta thu được $10$ gam muối ăn.
	}
\end{ex}

%%%%%============SA_2================%%%%%%
\begin{ex}%[6K3V3-1]
	Hỗn hợp A gồm bột sắt và bột đồng. Biết hỗn hợp A nặng $50$ gam. Dùng nam châm tách được toàn bộ bột sắt ra, cân lại phần còn lại được $18$ gam. Hỏi phần trăm khối lượng sắt trong hỗn hợp A là bao nhiêu? (Kết quả tính theo $\%$, không ghi đơn vị)
	\shortans{$64$}
	\loigiai{
		Sắt bị nam châm hút, đồng không bị nam châm hút.
		\\
		Sau khi dùng nam châm tách sắt, phần còn lại là đồng.
		\\
		Khối lượng bột đồng: $m_{\text{Cu}} = 18$ gam.
		\\
		Khối lượng bột sắt: $m_{\text{Fe}} = 50 - 18 = 32$ gam.
		\\
		Phần trăm khối lượng sắt trong hỗn hợp:
		\[ \%_{\text{Fe}} = \dfrac{m_{\text{Fe}}}{m_{\text{hỗn hợp}}} \times 100\% = \dfrac{32}{50} \times 100\% = 64\% \]
	}
\end{ex}

\Closesolutionfile{ansexh}
\Closesolutionfile{ansex}

%%%==============Phần bài tập tự luận==============%%%
\subsection{Bài tập tự luận}\textit{\large Thí sinh trả lời từ bài 1 đến bài 2}
\Opensolutionfile{ansbth}[Ans/LGBT-KTDK_KHTN_6_TachChat_MADE523]
\Opensolutionfile{ansbt}[Ans/AnsBT-KTDK_KHTN_6_TachChat_MADE523]

%%%%%============BT_1================%%%%%%
\begin{bt}%[6K3V4-1]
	Hỗn hợp B gồm muối ăn, bột sắt và cát. Em hãy trình bày các bước thí nghiệm để tách riêng từng chất ra khỏi hỗn hợp trên. Nêu rõ phương pháp và dụng cụ cần dùng ở mỗi bước.
	\loigiai{
		\textbf{Bước 1: Tách bột sắt bằng nam châm.}
		\\
		Dụng cụ: nam châm.
		\\
		Cách làm: Đưa nam châm vào hỗn hợp B, bột sắt có tính từ nên bị nam châm hút ra. Thu được bột sắt riêng. Hỗn hợp còn lại gồm muối ăn và cát.
		\\[6pt]
		\textbf{Bước 2: Hoà tan hỗn hợp muối và cát vào nước.}
		\\
		Dụng cụ: cốc thuỷ tinh, đũa khuấy, nước sạch.
		\\
		Cách làm: Cho hỗn hợp muối ăn và cát vào cốc nước, khuấy đều. Muối ăn tan hoàn toàn trong nước, cát không tan.
		\\[6pt]
		\textbf{Bước 3: Lọc để tách cát.}
		\\
		Dụng cụ: phễu lọc, giấy lọc, bình hứng.
		\\
		Cách làm: Đổ hỗn hợp qua phễu lọc có giấy lọc. Cát nằm lại trên giấy lọc (phần không tan). Nước muối chảy qua giấy lọc xuống bình hứng.
		\\[6pt]
		\textbf{Bước 4: Cô cạn nước muối để thu muối ăn.}
		\\
		Dụng cụ: bát sứ (chén sứ), đèn cồn, kiềng sắt.
		\\
		Cách làm: Đổ nước muối vào bát sứ, đun nóng trên ngọn lửa đèn cồn. Nước bay hơi dần, muối ăn kết tinh lại trong bát sứ. Thu được muối ăn tinh khiết.
		\\[6pt]
		\textbf{Kết luận:} Đã tách riêng được cả $3$ chất: bột sắt (bằng nam châm), cát (bằng lọc), muối ăn (bằng cô cạn).
	}
\end{bt}

%%%%%============BT_2================%%%%%%
\begin{bt}%[6K3V4-2]
	Ở vùng biển miền Trung Việt Nam, người dân thường làm muối bằng cách dẫn nước biển vào các ruộng muối rồi phơi dưới ánh nắng mặt trời.
	\begin{enumerate}[a)]
		\item Phương pháp làm muối từ nước biển thuộc phương pháp tách chất nào? Giải thích nguyên lý hoạt động.
		\item Nếu nước biển có nồng độ muối trung bình là $3{,}5\%$, tính khối lượng muối ăn thu được khi cô cạn $500$ kg nước biển. (Giả sử muối ăn là chất tan duy nhất trong nước biển)
		\item Tại sao ở Việt Nam, việc sản xuất muối thường tập trung vào mùa nắng (từ tháng $1$ đến tháng $8$)?
	\end{enumerate}
	\loigiai{
		\begin{enumerate}[a)]
			\item Phương pháp làm muối từ nước biển thuộc phương pháp \textbf{cô cạn} (bay hơi dung môi).
			\\
			Nguyên lý: Nước biển là hỗn hợp đồng nhất (dung dịch) gồm nước (dung môi) và muối ăn (chất tan). Khi phơi nắng, nước bay hơi dần nhờ nhiệt từ ánh sáng mặt trời. Muối ăn có nhiệt độ sôi rất cao ($1413\,^{\circ}\text{C}$) nên không bay hơi mà kết tinh lại trên ruộng muối.
			\item Khối lượng muối ăn thu được:
			\[ m_{\text{muối}} = \dfrac{C\% \times m_{\text{dung dịch}}}{100} = \dfrac{3{,}5 \times 500}{100} = 17{,}5 \text{ (kg)} \]
			Vậy từ $500$ kg nước biển, ta thu được $17{,}5$ kg muối ăn.
			\item Việc sản xuất muối tập trung vào mùa nắng vì:
			\begin{itemize}
				\item Mùa nắng có cường độ ánh sáng mặt trời lớn, nhiệt độ cao giúp nước bay hơi nhanh hơn, rút ngắn thời gian thu hoạch muối.
				\item Mùa nắng ít mưa, nếu mưa sẽ làm loãng nước muối trên ruộng, kéo dài thời gian cô cạn và giảm sản lượng muối.
				\item Số giờ nắng trong ngày dài hơn, tăng hiệu suất bay hơi nước.
			\end{itemize}
		\end{enumerate}
	}
\end{bt}

\Closesolutionfile{ansbt}
\Closesolutionfile{ansbth}

\begin{center}
 \rule[4pt]{2cm}{1pt}\,\large\bfseries Hết\,\rule[4pt]{2cm}{1pt}
\end{center}
\label{x}
\end{document}
